
\section{Relationen}


\dfn{Relation}{
  Seien $A$ und $B$ Mengen. Eine \textbf{Relation} $R$ zwischen
  $A$ und $B$ ist eine Teilmenge des kartesischen Produkts:
  \[
    R\subseteq A\times B.
  \]

  Das bedeutet, dass eine Relation bestimmte geordnete Paare
  aus $A\times B$ auswählt.

  Ist $A=B=M$, so heißt $R$ eine \textbf{Relation auf} $M$.
}

\nt{
  Eine Relation beschreibt, welche Elemente miteinander in
  Beziehung stehen. Sie muss nicht alle möglichen Paare enthalten.

  Beispielsweise seien $A$ die Menge der Studierenden und $B$
  die Menge der Professoren. Wir definieren
  \[
    (s,p)\in R
    \iff
    \text{die Person }s\text{ eine Vorlesung von }p\text{ hört}.
  \]
  Dann ist $R\subseteq A\times B$ eine Relation zwischen
  Studierenden und Professoren.

  Eine Relation kann auch auf einer einzigen Menge definiert sein.
  Beispielsweise beschreibt die Relation $\leq$ auf $\bbR$,
  welche reellen Zahlen kleiner oder gleich anderen sind.
}

\dfn{Eigenschaften von Relationen}{
  Sei $R$ eine Relation auf einer Menge $M$, also
  $R\subseteq M\times M$.

  \begin{enumerate}[label=\textbf{(\arabic*)}]
    \item \textbf{Reflexivität:}
    $R$ heißt reflexiv, falls jedes Element zu sich selbst
    in Relation steht:
    \[
      \forall a\in M:\ (a,a)\in R.
    \]

    \item \textbf{Symmetrie:}
    $R$ heißt symmetrisch, falls die Relation in beide Richtungen
    gilt:
    \[
      \forall a,b\in M:\ 
      (a,b)\in R\implies (b,a)\in R.
    \]

    \item \textbf{Antisymmetrie:}
    $R$ heißt antisymmetrisch, falls zwei Elemente, die in
    beiden Richtungen miteinander in Relation stehen, gleich
    sein müssen:
    \[
      \forall a,b\in M:\ 
      \bigl((a,b)\in R\land(b,a)\in R\bigr)
      \implies a=b.
    \]

    \item \textbf{Transitivität:}
    $R$ heißt transitiv, falls aus zwei aufeinanderfolgenden
    Relationen auch die direkte Relation folgt:
    \[
      \forall a,b,c\in M:\ 
      \bigl((a,b)\in R\land(b,c)\in R\bigr)
      \implies (a,c)\in R.
    \]

    \item \textbf{Totalität:}
    $R$ heißt total, falls je zwei Elemente miteinander
    vergleichbar sind:
    \[
      \forall a,b\in M:\ 
      (a,b)\in R\ \lor\ (b,a)\in R.
    \]
  \end{enumerate}
}

\nt{
  Die Begriffe lassen sich anhand der üblichen Relation $\leq$
  auf den reellen Zahlen verstehen.

  \begin{itemize}
    \item \textbf{Reflexiv:} Es gilt $a\leq a$ für jede Zahl $a$.

    \item \textbf{Nicht symmetrisch:} Aus $2\leq 5$ folgt nicht
    $5\leq 2$.

    \item \textbf{Antisymmetrisch:} Wenn $a\leq b$ und $b\leq a$,
    dann muss $a=b$ gelten.

    \item \textbf{Transitiv:} Aus $a\leq b$ und $b\leq c$ folgt
    $a\leq c$.

    \item \textbf{Total:} Für je zwei reelle Zahlen $a,b$ gilt
    stets $a\leq b$ oder $b\leq a$.
  \end{itemize}

  Besonders wichtig ist die Unterscheidung zwischen Symmetrie und
  Antisymmetrie. Symmetrie verlangt, dass die Relation auch in
  umgekehrter Richtung gilt. Antisymmetrie verlangt dagegen, dass
  zwei Elemente, die in beide Richtungen in Relation stehen,
  tatsächlich gleich sind.
}

\ex{Relationen auf einer Menge}{
  Sei $M=\{1,2,3\}$.

  Betrachte zunächst die Relation
  \[
    R_1=\{(1,1),(2,2),(3,3)\}.
  \]
  Jedes Element steht genau mit sich selbst in Relation.
  Deshalb ist $R_1$ reflexiv, symmetrisch, antisymmetrisch
  und transitiv.

  Betrachte nun
  \[
    R_2=\{(1,2),(2,1)\}.
  \]
  Diese Relation ist symmetrisch, weil sie sowohl $(1,2)$
  als auch $(2,1)$ enthält. Sie ist aber nicht reflexiv, da
  beispielsweise $(1,1)\notin R_2$ gilt.
}

\dfn{Äquivalenzrelation}{
  Eine Relation $R$ auf einer Menge $M$ heißt
  \textbf{Äquivalenzrelation}, falls sie die folgenden drei
  Eigenschaften erfüllt:
  \begin{enumerate}[label=\textbf{(\arabic*)}]
    \item Reflexivität,
    \item Symmetrie,
    \item Transitivität.
  \end{enumerate}
  Statt $(a,b)\in R$ schreibt man bei einer Äquivalenzrelation
  häufig $a\sim b$.
}

\nt{
  Eine Äquivalenzrelation beschreibt, wann zwei Elemente im
  betrachteten Zusammenhang als gleichartig angesehen werden.
  Die Elemente müssen dabei nicht identisch sein.

  Ein einfaches Beispiel ist die Relation auf $\bbZ$, bei der
  zwei ganze Zahlen genau dann miteinander in Relation stehen,
  wenn sie denselben Rest bei Division durch eine festgelegte
  natürliche Zahl haben.
}

\dfn{Ordnungsrelation}{
  Eine Relation $R$ auf einer Menge $M$ heißt
  \textbf{Ordnungsrelation} oder \textbf{Halbordnung}, falls sie
  reflexiv, antisymmetrisch und transitiv ist.

  Eine Ordnungsrelation heißt \textbf{Totalordnung}, falls
  zusätzlich je zwei Elemente vergleichbar sind, also
  \[
    \forall a,b\in M:\ 
    (a,b)\in R\lor(b,a)\in R.
  \]
}

\ex{Teilmengen als Ordnungsrelation}{
  Sei $M$ eine Menge und $\mathcal P(M)$ ihre Potenzmenge,
  also die Menge aller Teilmengen von $M$.

  Auf $\mathcal P(M)$ betrachten wir die Relation $\subseteq$.
  Für $A,B\in\mathcal P(M)$ gilt also
  \[
    A\subseteq B
    \iff
    \text{jedes Element von }A\text{ auch in }B\text{ liegt}.
  \]

  Diese Relation ist reflexiv, antisymmetrisch und transitiv:

  \begin{itemize}
    \item $A\subseteq A$ für jede Menge $A$.
    \item Aus $A\subseteq B$ und $B\subseteq A$ folgt $A=B$.
    \item Aus $A\subseteq B$ und $B\subseteq C$ folgt
    $A\subseteq C$.
  \end{itemize}

  Daher ist $\subseteq$ eine Ordnungsrelation auf $\mathcal P(M)$.

  Sie ist im Allgemeinen keine Totalordnung: Zwei Mengen müssen
  nicht ineinander enthalten sein. Beispielsweise sind
  $\{1\}$ und $\{2\}$ nicht miteinander vergleichbar.
}

\nt{
  Eine Halbordnung verlangt also nicht, dass sich alle Elemente
  vergleichen lassen. Bei einer Totalordnung ist dies dagegen
  vorgeschrieben. Die Relation $\leq$ auf $\bbR$ ist eine
  Totalordnung, während die Teilmengenrelation $\subseteq$
  im Allgemeinen nur eine Halbordnung ist.
}

\dfn{Kongruenz modulo $m$}{
  Sei $m\in\bbN$ mit $m\geq 1$. Für $a,b\in\bbZ$ definieren wir
  \[
    a\equiv b\pmod m
    \iff
    m\mid(a-b).
  \]
  Das bedeutet, dass $a-b$ durch $m$ teilbar ist.

  Gleichwertig dazu haben $a$ und $b$ bei Division durch $m$
  denselben Rest.
}

\mlemma{Die Kongruenz modulo $m$ ist eine Äquivalenzrelation}{
  Für jedes $m\in\bbN$ mit $m\geq1$ ist die Relation
  $\equiv\pmod m$ auf $\bbZ$ reflexiv, symmetrisch und transitiv.
}

\begin{myproof}
  \textbf{Reflexivität:}

  Für jedes $a\in\bbZ$ gilt
  \[
    a-a=0.
  \]
  Da $m\mid0$ gilt, folgt
  \[
    a\equiv a\pmod m.
  \]

  \textbf{Symmetrie:}

  Seien $a,b\in\bbZ$ und es gelte
  \[
    a\equiv b\pmod m.
  \]
  Dann ist $m\mid(a-b)$. Es gibt also ein $s\in\bbZ$ mit
  \[
    a-b=sm.
  \]
  Durch Multiplikation mit $-1$ erhalten wir
  \[
    b-a=(-s)m.
  \]
  Somit gilt $m\mid(b-a)$ und daher
  \[
    b\equiv a\pmod m.
  \]

  \textbf{Transitivität:}

  Seien $a,b,c\in\bbZ$ und es gelte
  \[
    a\equiv b\pmod m
    \qquad\text{und}\qquad
    b\equiv c\pmod m.
  \]
  Dann existieren $s,t\in\bbZ$ mit
  \[
    a-b=sm,\qquad b-c=tm.
  \]
  Durch Addition folgt
  \[
    a-c=(a-b)+(b-c)=(s+t)m.
  \]
  Also ist $m\mid(a-c)$ und somit
  \[
    a\equiv c\pmod m.
  \]

  Damit ist die Relation reflexiv, symmetrisch und transitiv,
  also eine Äquivalenzrelation.
\end{myproof}

\ex{Kongruenz modulo $3$}{
  Zwei ganze Zahlen sind genau dann kongruent modulo $3$,
  wenn sie bei Division durch $3$ denselben Rest haben.

  Beispielsweise gilt
  \[
    8\equiv 2\pmod 3,
  \]
  denn $8-2=6$ ist durch $3$ teilbar. Ebenso gilt
  \[
    5\equiv 2\pmod 3,
    \qquad
    8\equiv 5\pmod 3.
  \]

  Dagegen gilt
  \[
    8\not\equiv 4\pmod 3,
  \]
  denn $8-4=4$ ist nicht durch $3$ teilbar.
}

\dfn{Äquivalenzklasse}{
  Sei $\sim$ eine Äquivalenzrelation auf einer Menge $M$.
  Für $a\in M$ heißt die Menge
  \[
    [a]:=\{x\in M\mid x\sim a\}
  \]
  die \textbf{Äquivalenzklasse} von $a$.

  Sie besteht aus allen Elementen von $M$, die zu $a$
  in Relation stehen.
}

\nt{
  Die Schreibweise $[a]$ bezeichnet nicht nur das einzelne
  Element $a$, sondern eine ganze Menge von Elementen.

  Da eine Äquivalenzrelation reflexiv ist, gilt immer
  $a\sim a$. Somit liegt $a$ selbst in seiner Äquivalenzklasse:
  \[
    a\in[a].
  \]

  Bei der Kongruenz modulo $3$ ist zum Beispiel
  \[
    [0]=\{\dots,-6,-3,0,3,6,9,\dots\}.
  \]
  Diese Klasse enthält alle durch $3$ teilbaren ganzen Zahlen.
  Entsprechend enthält $[1]$ alle Zahlen mit Rest $1$ und
  $[2]$ alle Zahlen mit Rest $2$.
}

\thm{Eigenschaften von Äquivalenzklassen}{
  Sei $\sim$ eine Äquivalenzrelation auf einer Menge $M$.
  Dann gilt:

  \begin{enumerate}[label=\textbf{(\alph*)}]
    \item Jede Äquivalenzklasse ist nicht leer:
    \[
      [a]\neq\emptyset
      \qquad\text{für alle }a\in M.
    \]

    \item Die Vereinigung aller Äquivalenzklassen ist ganz $M$:
    \[
      \bigcup_{a\in M}[a]=M.
    \]

    \item Zwei Äquivalenzklassen sind entweder gleich oder
    disjunkt:
    \[
      [a]=[b]
      \quad\text{oder}\quad
      [a]\cap[b]=\emptyset
      \qquad\text{für alle }a,b\in M.
    \]
  \end{enumerate}
}

\begin{myproof}
  \textbf{(a) Nichtleere Äquivalenzklassen:}

  Sei $a\in M$ beliebig. Wegen der Reflexivität gilt $a\sim a$.
  Daher ist $a\in[a]$, also
  \[
    [a]\neq\emptyset.
  \]

  \textbf{(b) Vereinigung aller Äquivalenzklassen:}

  Sei $x\in M$. Wegen $x\sim x$ gilt $x\in[x]$. Somit liegt
  jedes Element von $M$ in mindestens einer Äquivalenzklasse:
  \[
    M\subseteq\bigcup_{a\in M}[a].
  \]
  Umgekehrt ist jede Äquivalenzklasse eine Teilmenge von $M$.
  Daher gilt auch
  \[
    \bigcup_{a\in M}[a]\subseteq M.
  \]
  Insgesamt folgt
  \[
    \bigcup_{a\in M}[a]=M.
  \]

  \textbf{(c) Gleichheit oder Disjunktheit:}

  Seien $a,b\in M$. Wir nehmen an, dass
  \[
    [a]\cap[b]\neq\emptyset.
  \]
  Dann gibt es ein Element $x\in M$ mit
  \[
    x\in[a]\quad\text{und}\quad x\in[b].
  \]
  Nach Definition der Äquivalenzklassen bedeutet dies
  \[
    x\sim a\quad\text{und}\quad x\sim b.
  \]
  Wegen der Symmetrie folgt $a\sim x$. Zusammen mit $x\sim b$
  ergibt die Transitivität
  \[
    a\sim b.
  \]

  Nun zeigen wir $[a]\subseteq[b]$. Sei dazu $y\in[a]$.
  Dann gilt $y\sim a$. Aus $a\sim b$ folgt mit der
  Transitivität
  \[
    y\sim a\quad\text{und}\quad a\sim b
    \implies y\sim b.
  \]
  Also ist $y\in[b]$ und damit $[a]\subseteq[b]$.

  Analog zeigt man $[b]\subseteq[a]$. Somit gilt
  $[a]=[b]$.

  Wir haben gezeigt: Sobald zwei Äquivalenzklassen ein
  gemeinsames Element besitzen, sind sie gleich. Andernfalls
  ist ihr Schnitt leer.
\end{myproof}

\nt{
  Der Satz bedeutet, dass eine Äquivalenzrelation die Menge $M$
  in disjunkte Klassen aufteilt. Jedes Element gehört zu genau
  einer solchen Klasse.

  Man kann sich das so vorstellen: Zunächst werden Elemente,
  die miteinander in Relation stehen, zu Gruppen zusammengefasst.
  Aufgrund der Eigenschaften einer Äquivalenzrelation überlappen
  sich diese Gruppen nicht – sie sind entweder identisch oder
  haben kein gemeinsames Element.

  Eine Teilmenge von $M$, die aus jeder Äquivalenzklasse genau
  ein Element enthält, heißt \textbf{Repräsentantensystem}.
  Ein solches System muss im Allgemeinen nicht eindeutig sein.
}

\ex{Äquivalenzklassen der Gleichheitsrelation}{
  Sei $M$ eine Menge und definiere
  \[
    a\sim b\iff a=b.
  \]
  Dies ist eine Äquivalenzrelation auf $M$.

  Für jedes $a\in M$ ist die Äquivalenzklasse
  \[
    [a]=\{a\}.
  \]
  Jede Klasse enthält also genau ein Element.
}

\ex{Äquivalenzklassen modulo $m$}{
  Betrachte die Kongruenzrelation modulo $m$ auf $\bbZ$.
  Zwei ganze Zahlen gehören genau dann zur selben
  Äquivalenzklasse, wenn sie denselben Rest bei Division
  durch $m$ haben.

  Es gibt genau $m$ verschiedene Äquivalenzklassen. Man kann
  sie durch die Vertreter $0,1,\dots,m-1$ beschreiben:
  \[
    [0],[1],\dots,[m-1].
  \]

  Beispielsweise gilt für $m=3$:
  \begin{align*}
    [0] &= \{\dots,-6,-3,0,3,6,\dots\},\\
    [1] &= \{\dots,-5,-2,1,4,7,\dots\},\\
    [2] &= \{\dots,-4,-1,2,5,8,\dots\}.
  \end{align*}

  Jede ganze Zahl liegt in genau einer dieser Klassen, je nachdem,
  ob ihr Rest bei Division durch $3$ gleich $0$, $1$ oder $2$ ist.
}
